现在进一步减弱对 Dirichlet 问题的解的条件: 将问题陈述为变分形式. 这与 Chapter~\ref{chap:conservation-laws-shocks} 中讨论守恒律弱解时使用的过程相同. 第一步是利用分部积分, 从微分算子 $L$ 构造双线性型. 设 $\phi\in\mathcal{D}(\Omega)$ 且 $u\in H^{2k}(\Omega)$, 则
\begin{equation}\label{eq:dirichlet-variational-integration-by-parts}
\begin{aligned}
\int_\Omega\phi Lu\,\mathrm{d}\bm{x}
&=\sum_{0\leq|\sigma|,|\gamma|\leq k}(-1)^{|\sigma|}\int_\Omega\phi D^\sigma(a_{\sigma\gamma}(\bm{x})D^\gamma u)\,\mathrm{d}\bm{x}\\
&=\sum_{0\leq|\sigma|,|\gamma|\leq k}\int_\Omega a_{\sigma\gamma}(\bm{x})D^\gamma uD^\sigma\phi\,\mathrm{d}\bm{x}.
\end{aligned}
\end{equation}


据此定义
\begin{equation}\label{eq:elliptic-dirichlet-bilinear-form-definition}
B[v,u]:=\sum_{0\leq|\sigma|,|\gamma|\leq k}\int_\Omega a_{\sigma\gamma}(\bm{x})D^\gamma uD^\sigma v\,\mathrm{d}\bm{x}
\end{equation}
为与椭圆偏微分算子 $L$ 相联系的\textbf{双线性型}. 注意, 即使 $u$ 和 $v$ 仅属于 $H^k(\Omega)$, $B[v,u]$ 也有定义. 据此, 给出 Dirichlet 问题的又一种解的定义.

\begin{Definition}
\label{def:weak-dirichlet-solution}
设 $\Omega\subset\mathbb{R}^n$ 为有界区域, 给定 $f\in H^{-k}(\Omega)$. 若函数\[u\in H_0^k(\Omega)\]对每个 $v\in H_0^k(\Omega)$ 都满足
\begin{equation}\label{eq:weak-dirichlet-variational-equation}
B[v,u]=f(v)
\end{equation}
则称其为 Dirichlet 问题的\textbf{弱解}.
\end{Definition}

\begin{Remark}
\label{rem:complex-dirichlet-sesquilinear-extension}
可令
\begin{equation}\label{eq:complex-elliptic-dirichlet-sesquilinear-form}
B[v,u]:=\sum_{0\leq|\sigma|,|\gamma|\leq k}\int_\Omega a_{\sigma\gamma}(\bm{x})D^\gamma u\overline{D^\sigma v}\,\mathrm{d}\bm{x}.
\end{equation}
从而将实 Hilbert 空间 $H_0^k(\Omega)$ 上的双线性型 $B$ 延拓为复 Hilbert 空间 (仍记为 $H_0^k(\Omega)$) 上的\emph{半双线性型}. 对复值函数与实值函数的区别不作十分严格的区分, 相应空间及双线性型均使用同一记号. 由于主要关心偏微分方程的实解, 只有不得不使用复空间时才转向复空间, 例如使用 Fourier 变换或讨论谱理论时.
\end{Remark}

注意, 利用 \eqref{eq:dirichlet-variational-integration-by-parts} 中的计算 (不过这次以函数 $v\in H_0^k(\Omega)$ 取代 $\phi\in\mathcal{D}(\Omega)$), 可见 Dirichlet 问题的任意强解自然也是弱解. 然而, 由于对弱解的光滑性要求远低于强解, 证明弱解存在会容易得多. 做到这一点后, 将能证明: 若 $\Omega$, $f$ 和系数 $a_{\sigma\gamma}$ 足够 ``好\textquotesingle\textquotesingle, 则弱解事实上是强解或经典解.
